\documentclass{article}

\usepackage{tocloft}
\usepackage{tabularx}
\usepackage{graphicx}
\usepackage{array}
\usepackage{tabu}
\usepackage[table]{xcolor}

%creazione dei colori custom
\definecolor{tableGreen}{RGB}{25, 200, 60}
\definecolor{tableRed}{RGB}{220 ,15, 5}
\definecolor{tableBlue}{RGB}{0 ,20, 210}
\definecolor{tableCyan}{RGB}{0 ,175, 185}

%macro per andare a capo nelle tabelle
\newcommand{\n}{
   \\\hline 
}

%macro per generare le tabelle
\newcommand{\monkeytable}[2]{

\begin{center}
\begin{table}[h!]
\setlength{\extrarowheight}{5pt} %padding
\begin{tabularx}{\textwidth}
#2
\n %necessaria per chiudere la tabella
\end{tabularx}
\end{table}
\label{tab:monkeytable}
\end{center}
}




%inizio documento vero e proprio
\title{Codemonkey}

\begin{document}

\begin{figure}[!h]
   \centering
   \includegraphics[width=0.5\linewidth]{images/icon}
   \begin{titlepage}
      \huge Codemonkey
   \end{titlepage}
   \label{fig:nome-etichetta}
\end{figure}

\pagebreak
\tableofcontents
\pagebreak
\section {Abstract}

L'obiettivo è sviluppare una piattaforma web che espone programmatori a realtà lavorative che ne hanno bisogno.
Sia i programmatori che le aziende si registrano al sito, e sono interfacciati al servizio Codemonkey attraverso pannelli diversi.
Un programmatore registrato inserisce i propri contatti, una breve descrizione e un elenco degli strumenti software nei quali ha competenze.
Un'azienda registrata può cercare programmatori filtrandoli per gli strumenti software dei quali ha esigenza. Quando l'azienda sceglie un programmatore, a questo arriva una notifica e può decidere di accettare il lavoro.
A lavoro concluso, l'azienda valuta il programmatore con un sistema di quotazione.
Il profilo del programmatore è aggiornato con i suoi lavori recenti.
\pagebreak
\section {\Large Raccolta dei requisiti}
\begin{itemize}
\large
\item Gli utenti del servizio Codemonkey si registrano da una pagina e indicano se intendono registrarsi come Codmonkey o Clienti
\item La Codmonkey potrá accedere, modificare i suoi dati e accettare eventuali lavori dopo aver eseguito l'Autenticazione
\item Un Cliente potrá accedere, modificare i suoi dati e gestire le Collaborazioni dopo aver eseguito il l'Autenticazione
\item Gli Utenti possono sfogliare il sito. Solo i clienti possono mandare richieste di Lavoro alle Codmonkey
\item Le Codmonkey possono accettare o rifiutare i lavori
\item I clienti possono valutare la Codmonkey solo dopo aver terminato la collaborazione e la valutazione potrá ache opzionalmente contenere una recensione
\item Si potranno cercare Codmonkey utlizzando dei filti specifici
\end{itemize}

\pagebreak
\section {Vocabolario}

\begin{table}[ht]
\begin{tabularx}{\textwidth}{|X|X|X|}
\hline
Voce & Definizione & Sinonimo \\
\hline
Programmatore & Utente registrato che fornisce uno o piú servizi alle aziende  & \\
\hline
Azienda & Azienda o un semplice privato interessato a utilizzare uno o piú servizi offerti da un programmatore & \\
\hline
Credenziali  & Metodo di accesso al servizio, basato su username e password & \\
\hline
Account & Insieme di Credenziali e informazioni che identifica una Azienda/Programmatore & \\
\hline
Utente & \begin{tabular}{@{}l@{}}Utilizzatore del servizio \\(Sia Azienda che \\ Programmatore)\end{tabular} & \\
\hline
Username & \begin{tabular}{@{}l@{}} Stringa alfanumerica \\ che identifica il nome \\dell'utente che sta \\ tentando l'accesso \end{tabular} & Identificativo\\
\hline
Password & \begin{tabular}{@{}l@{}}Stringa alfanumerica\\ generata da un utente\\ del servizio\end{tabular} & \\
\hline
Autenticazione & \begin{tabular}{@{}l@{}}Meccanismo di accesso \\ con nome alla \\piattaforma\end{tabular} &  Log in\\
\hline
Registrazione & Funzione di iscrizione alla piattaforma & \\
\hline
\end{tabularx}
\end{table}
\pagebreak
\section {Tabelle dei Requisiti}

\newcounter{greenC}

\newcommand{\tableGreen}{%onde evitare caos il colore della tabella viene dichiarato qui
    \\
    \rowcolor{tableGreen!15}
    \hline
    \stepcounter{greenC}
    R\thegreenC F
}

\newcommand{\ntableGreen}{
    \\
    \rowcolor{tableGreen!5}
    \hline
    \stepcounter{greenC}
    R\thegreenC F
}

\monkeytable{3} {
{|>{\arraybackslash}m{1.5cm}|>{\centering\arraybackslash}X|}

\hline
\rowcolor{tableGreen!70}
\multicolumn{2}{|c|}{\textbf{Requisiti Funzionali}}
\n \rowcolor{tableGreen!50} \textbf{ID} & \textbf{Requisito}
\ntableGreen    & un utente qualsiasi può sfogliare il sito e vedere tutti i programmatori
\tableGreen     & il sito deve fornire la possibilità di inserire dei filtri per cercare programmatori con specifiche caratteristiche
\ntableGreen    & è presente una sezione nella quale un utente o una azienda possono autenticarsi
\tableGreen     & è presente una sezione nella quale un utente o una azienda possono registrarsi
\ntableGreen    & una azienda potrà contattare un programmatore solo previa autenicazione
\tableGreen     & un programmatore può accedere previa autenticazione alla pagina del suo profilo
\ntableGreen    & possono essere effettuate modifiche delle informazioni del programmatore dalla pagina principale
\tableGreen     & possono essere accettati i lavori e mandati i preventivi sempre da essa
\ntableGreen    & viene fornita la valutazione di ogni programmatore anche filtrata secondo specifici linguaggi
\tableGreen     & viene fornita la possibilità di vedere se un programmatore sta lavorando ad un progetto o meno
}
\pagebreak
\newcounter{yellowC}

\newcommand{\Yellow}{%onde evitare caos il colore della tabella viene dichiarato qui
    \\
    \rowcolor{tableYellow!15}
    \hline
    \stepcounter{yellowC}
    R\theyellowC NF
}

\newcommand{\nYellow}{
    \\
    \rowcolor{tableYellow!5}
    \hline
    \stepcounter{yellowC}
    R\theyellowC NF
}
\monkeytable{3} {
{|>{\arraybackslash}m{1.5cm}|>{\arraybackslash}X|}

\hline
\rowcolor{tableYellow!70}
\multicolumn{2}{|c|}{\textbf{Requisiti Non Funzionali}}
\n \rowcolor{tableYellow!50} \textbf{ID} & \centering\textbf{Requisito} \endline
\rowcolor{tableYellow!15}
\hline
\stepcounter{yellowC}
R\theyellowC NF 
            & Il sito deve essere facile da navigare
\nYellow    & Deve essere tracciata l'attivitá dei vari amministratori
\Yellow     & Viene fornita la Valutazione Generale di ogni Codmonkey
\nYellow    & Viene fornita una Valutazione Filtrata del programmatore in base ai filtri impostati nella ricerca del Cliente
\Yellow     & Deve essere possibile vedere a qunti progetti la Codmonkey sta lavorando

}
\pagebreak
\newcounter{reqDCounter}%counter

\newcommand{\nReqD}{
    \n
    \stepcounter{reqDCounter}
    R\thereqDCounter D
}

\begin{center}
    \rowcolors{2}{tableBlue!4}{tableBlue!8}%colori alternati
    \begin{tabularx}{\textwidth}
        {|c|X|}
        \hline\rowcolor{tableBlue!50}
        \multicolumn{2}{|c|}{\Large\textbf{Requisiti di Dominio}}
        \n \rowcolor{tableBlue!25} \large\textbf{ID}
               & \large\textbf{Requisito}
        \nReqD & Non é necessaria l'autenticazione per eseguire una rcerca di Codmonkey
        \nReqD & Gli amministratori del sistema devono essere creati prima dell'avvio del sistema e non se ne possono aggiungere durante l'esecuzione
        \nReqD & Non sará possibile regisrtarsi si come codmonkey che come Cliente
        \nReqD & Un Utente Registrato in qualsiasi momento puó richiedere l'eliminazione del suo Account
        \n
    \end{tabularx}
\end{center}
\pagebreak
\section {Scenari}


\newcommand{\tableCyan}{%onde evitare caos il colore della tabella viene dichiarato qui
    \\
    \hline 
    \rowcolor{tableCyan!10}
    \cellcolor{tableCyan!22}
}

\newcommand{\ntableCyan}{
    \\
    \hline
    \rowcolor{tableCyan!5}
    \cellcolor{tableCyan!12}
}


\monkeytable{3} {
{|>{\arraybackslash }m{3cm}|>{\arraybackslash}X|}

\hline \rowcolor{tableCyan!37} \centering\textbf{Titolo} & \centering\textbf{...}\endline
\hline 
\rowcolor{tableCyan!10}
\cellcolor{tableCyan!22}

               Descrizione &test
\ntableCyan    Attori &
\tableCyan     Relazioni &
\ntableCyan    Precondizioni &
\tableCyan     Postcondizioni &
\ntableCyan    Scenario Principale &
\tableCyan     Scenari Alternativi &
\ntableCyan    Requisiti NF &
\tableCyan     Punti Aperti &
}










    % NO    % NO    % NO    % NO    % NO
% NO    % NO    % NO    % NO    % NO

%Si usa per il copia e incolla
%Non  riempitela scimmieeeee                                       monke

\monkeytable{3} {
{|>{\arraybackslash }m{3cm}|>{\arraybackslash}X|}

\hline \rowcolor{tableCyan!37} \centering\textbf{Titolo} & \centering\textbf{...}\endline
\hline 
\rowcolor{tableCyan!10}
\cellcolor{tableCyan!22}

               Descrizione &test
\ntableCyan    Attori &
\tableCyan     Relazioni &
\ntableCyan    Precondizioni &
\tableCyan     Postcondizioni &
\ntableCyan    Scenario Principale &
\tableCyan     Scenari Alternativi &
\ntableCyan    Requisiti NF &
\tableCyan     Punti Aperti &
}

    % NO    % NO    % NO    % NO    % NO
% NO    % NO    % NO    % NO    % NO
                                                                                                                                         %un simpatico easter egg

\end{document}
